% ============================================================
% TITLE PAGE — SGCV-IA — Requirements Engineering (Springer)
%
% Archivo independiente con información de autoría y declaraciones.
% Mantener separado del manuscrito anonimizado cuando el proceso de
% revisión de la revista así lo requiera.
% ============================================================

\documentclass[12pt]{article}

\usepackage[margin=2.5cm]{geometry}
\usepackage[T1]{fontenc}
\usepackage[utf8]{inputenc}
\usepackage{hyperref}

\begin{document}

% ============================================================
% TITLE
% ============================================================

\begin{center}
{\Large\bfseries
Explainability as a non-functional requirement in an AI-assisted
veterinary clinical system: a case study
}
\end{center}

\vspace{1em}

% ============================================================
% AUTHORS
% ============================================================

\section*{Authors}

\begin{enumerate}

  \item Alberto Jeanpool Marcillo Ponce$^{1,\ast}$\\
  ORCID: \href{https://orcid.org/0009-0003-7886-3957}
  {0009-0003-7886-3957}\\
  \texttt{amarcillop@uteq.edu.ec}

  \item Robyn Willian Amagua Sacón$^{1}$\\
  ORCID: \href{https://orcid.org/0009-0006-7363-4276}
  {0009-0006-7363-4276}\\
  \texttt{ramaguas@uteq.edu.ec}

  \item Anthony Alfredo Vera Gómez$^{1}$\\
  ORCID: \href{https://orcid.org/0009-0007-8198-6199}
  {0009-0007-8198-6199}\\
  \texttt{averag10@uteq.edu.ec}

  \item Jhon Alexander Mesías Quijije$^{1}$\\
  ORCID: \href{https://orcid.org/0009-0001-6929-7514}
  {0009-0001-6929-7514}\\
  \texttt{jmesiasq@uteq.edu.ec}

  \item Carlos Daniel Barrionuevo Fuentes$^{1}$\\
  ORCID: \href{https://orcid.org/0009-0003-3415-3208}
  {0009-0003-3415-3208}\\
  \texttt{cbarrionuevof@uteq.edu.ec}

  \item Gleiston Ciceron Guerrero Ulloa$^{1}$ (academic supervisor)\\
  ORCID: \href{https://orcid.org/0000-0001-5990-2357}
  {0000-0001-5990-2357}\\
  \texttt{gguerrero@uteq.edu.ec}

\end{enumerate}

\noindent
$^{1}$Facultad de Ciencias de la Computación,
Universidad Técnica Estatal de Quevedo (UTEQ),
Quevedo, Ecuador.

\vspace{0.5em}

\noindent
$^{\ast}$Corresponding author:
Alberto Jeanpool Marcillo Ponce,
\texttt{amarcillop@uteq.edu.ec}



% ============================================================
% ABSTRACT
% ============================================================

\section*{Abstract}

\noindent\textbf{Context.} AI-assisted diagnostic modules in veterinary clinical-management systems frequently lack explainability operationalized as a verifiable requirement.

\noindent\textbf{Objective.} To determine which explainability requirements, operationalized as verifiable non-functional requirements (NFRs), are necessary and sufficient for the different user profiles (technical and non-technical) of an AI-assisted veterinary clinical system.

\noindent\textbf{Method.} Applied single-case study grounded in Requirements Engineering. Requirements were elicited through semi-structured interviews and direct observation, analyzed qualitatively, operationalized as functional and non-functional requirements, and linked through requirements traceability.

\noindent\textbf{Results.} Open coding of 16 semi-structured interviews yielded 167 open codes, consolidated into 50 axial codes across 7 categories; closed-coding saturation was reached (1.33 new axial codes on average in the final three interviews, 2.67\% of the cumulative total, below the 5\% threshold). Acceptance of AI-assisted diagnostic suggestions was consistently conditional: participants required the veterinarian's ability to review and correct suggestions (4/16) and grounding in validated literature and data-protection safeguards (6/16), with general moderate-to-high acceptance (11/16) and AI treated as decision support rather than a substitute for professional judgement. These conditions were operationalized as non-functional requirement RNF-18, traced to the AI-assisted diagnostic use case CU-06.

\noindent\textbf{Conclusions.} Explainability for AI-assisted veterinary diagnostic support is best operationalized around contestability and provenance of suggestions rather than technical model interpretability alone, and can be linked, through requirements traceability, to concrete verifiable requirements. Broader validation across additional cases and a second independent coding round remain as future work.

\textit{Note: this abstract must remain byte-for-byte identical to the abstract in the anonymized manuscript (\texttt{manuscrito\_final.tex}). If either is edited, update both.}

% ============================================================
% KEYWORDS
% ============================================================

\section*{Keywords}

Explainability; Non-functional requirements; AI-assisted diagnosis;
Veterinary informatics; User studies

\vspace{1.5em}

% ============================================================
% DECLARATIONS
% ============================================================

\section*{Declarations}

% ------------------------------------------------------------
% FUNDING
% ------------------------------------------------------------

\subsection*{Funding}

The authors received no specific funding for this work.

% ------------------------------------------------------------
% COMPETING INTERESTS
% ------------------------------------------------------------

\subsection*{Competing Interests}

The authors declare no competing interests.

% ------------------------------------------------------------
% ETHICS
% ------------------------------------------------------------

\subsection*{Ethics approval and consent to participate}

Participation in the requirements-elicitation activities was voluntary
and based on written informed consent. Participant information is
handled using anonymized identifiers, and identifying materials are
maintained separately in the restricted evidence area in accordance
with the project's data-protection procedures and Ecuador's
\textit{Ley Orgánica de Protección de Datos Personales}
(LOPDP).

\par
No individual ethics-committee approval identifier was issued for
this project. This is consistent with the course-level ethical
review process applied uniformly across all course teams for the
Ingeniería de Requerimientos (ISR-401) Proyecto Fin de Curso at
Universidad Técnica Estatal de Quevedo, rather than an omission
specific to this study. \textit{Pending confirmation before
submission:} the authors should verify with the course instructor
whether a general/institutional approval reference exists that
should be cited here on behalf of the course as a whole (see
\texttt{08\_Etica/Aval\_Institucional.pdf}).

% ------------------------------------------------------------
% CONSENT FOR PUBLICATION
% ------------------------------------------------------------

\subsection*{Consent for publication}

\textit{Pending verification before submission.}
The final statement must reflect the exact scope of the signed consent
documents regarding the use and publication of anonymized participant
information.

% ------------------------------------------------------------
% DATA / MATERIALS / CODE
% ------------------------------------------------------------

\subsection*{Data, Materials and/or Code availability}

The research materials supporting the requirements-engineering phase
include anonymized interview transcripts, thematic-coding materials,
requirements artefacts, traceability information, and related project
documentation. Potentially identifying materials, including original
consent documents and multimedia evidence, are maintained separately
in the restricted evidence area and are not intended for public
dissemination.

\par
The anonymized replication package is deposited on Zenodo:
\href{https://doi.org/10.5281/zenodo.22238486}{https://doi.org/10.5281/zenodo.22238486}.
\textit{Pending verification before submission: this DOI currently
points to the automatically generated deposit under an MIT license;
it must be replaced with the DOI of the curated CC~BY~4.0 dataset
package once that deposit is finalized (see \texttt{07\_Datos/registro\_deposito.md}).}
The protocol for the Enfoque 3 empirical component is registered on
OSF: \href{https://osf.io/r5p8d/}{https://osf.io/r5p8d/} (registered
1 August 2026).

% ------------------------------------------------------------
% AUTHOR CONTRIBUTIONS
% ------------------------------------------------------------

\subsection*{Author contributions}

\textit{Partial draft based on verified repository artefacts; must be
reviewed and completed by all authors before submission, preferably
following the CRediT taxonomy.}

\textbf{Marcillo Ponce, A.J.}: thematic coding (open coding of
interviews P09--P16; axial consolidation), thematic-saturation
analysis, traceability-matrix verification, manuscript writing
(original draft and revisions), project administration for this
submission.
\textbf{Amagua Sacón, R.W.}: traceability-matrix corrections and
verification, open-coding saturation analysis script and figure.
\textbf{Vera Gómez, A.A.}: Enfoque 3 empirical-component design,
OSF protocol registration, validation-instrument design, confirmation
of protocol deviations.
\textbf{Mesías Quijije, J.A.}: \textit{pending -- to be completed by
the author.}
\textbf{Barrionuevo Fuentes, C.D.}: \textit{pending -- to be completed
by the author.}
\textbf{Guerrero Ulloa, G.C.}: academic supervision.

All authors read and approved the final manuscript.

% ------------------------------------------------------------
% AI-ASSISTED TECHNOLOGIES
% ------------------------------------------------------------

\subsection*{Use of AI-assisted technologies}

During preparation and revision of this manuscript, an AI-assisted
large language model (Claude, Anthropic) was used interactively,
with human oversight, exclusively as a support tool for verification
and language improvement. Its use was limited to: (i) checking the
consistency, clarity, grammar, spelling, and wording of text previously
prepared by the authors; (ii) verifying bibliographic and other factual
data against the information and sources used by the authors; and
(iii) suggesting improvements to the writing and overall clarity of
the manuscript. The AI tool was not used to generate original or
substantive content, including research data, thematic codes,
categories, analysis, results, methodology, interpretations,
conclusions, tables, figures, or research decisions. All suggestions
were reviewed by the authors before inclusion. The authors remain
fully responsible for the accuracy, provenance, and scientific
integrity of all content and for the final version of the manuscript.

% ------------------------------------------------------------
% ACKNOWLEDGEMENTS
% ------------------------------------------------------------

\subsection*{Acknowledgements}

The authors thank the veterinary professionals, administrative
personnel, and other participants who contributed to the
requirements-elicitation and validation activities of the SGCV-IA
project. Participant identities are not disclosed in this manuscript
in accordance with the project's anonymization and data-protection
procedures.

\end{document}
